\SourceLeaf{260}%
Introducendo has mensuras curvaturae in serie pro \(\sigma\), obtinemus expressionem sequentem, usque ad quantitates sexti ordinis (excl.) exactam:
\[
\begin{aligned}
\sigma=\tfrac{1}{2}ac\sin B\{&1+\tfrac{1}{120}\alpha(4pp-2qq+3qq'+3q'q')\\
&+\tfrac{1}{120}\beta(3pp-6qq+6qq'+3q'q')\\
&+\tfrac{1}{120}\gamma(3pp-2qq+qq'+4q'q')\}
\end{aligned}
\]
Praecisio eadem manebit, si pro \(p,q,q'\) substituimus \(c\sin B,c\cos B,c\cos B-a\), quo pacto prodit
\[
\begin{aligned}
\relax[8]\quad\sigma=\tfrac{1}{2}ac\sin B\{&1+\tfrac{1}{120}\alpha(3aa+4cc-9ac\cos B)\\
&+\tfrac{1}{120}\beta(3aa+3cc-12ac\cos B)\\
&+\tfrac{1}{120}\gamma(4aa+3cc-9ac\cos B)\}
\end{aligned}
\]
Quum ex hac aequatione omnia, quae ad lineam \(AD\) normaliter ad \(BC\) ductam referuntur, evanuerint, etiam puncta \(A,B,C\) cum correlatis inter se permutare licebit, quapropter erit eadem praecisione
\[
\begin{aligned}
\relax[9]\quad\sigma=\tfrac{1}{2}bc\sin A\{&1+\tfrac{1}{120}\alpha(3bb+3cc-12bc\cos A)\\
&+\tfrac{1}{120}\beta(3bb+4cc-9bc\cos A)\\
&+\tfrac{1}{120}\gamma(4bb+3cc-9bc\cos A)\}
\end{aligned}
\]
\[
\begin{aligned}
\relax[10]\quad\sigma=\tfrac{1}{2}ab\sin C\{&1+\tfrac{1}{120}\alpha(3aa+4bb-9ab\cos C)\\
&+\tfrac{1}{120}\beta(4aa+3bb-9ab\cos C)\\
&+\tfrac{1}{120}\gamma(3aa+3bb-12ab\cos C)\}
\end{aligned}
\]
\Article{26}
Magnam utilitatem affert consideratio trianguli plani rectilinei, cuius latera aequalia sunt ipsis \(a,b,c\); anguli illius trianguli, quos per \(A^*,B^*,C^*\) designabimus, different ab angulis trianguli in superficie curva, puta ab \(A,B,C\), quantitatibus secundi ordinis, operaeque pretium erit, has differentias accurate evolvere. Calculorum autem prolixiorum quam difficiliorum, primaria momenta apposuisse sufficiet.

Mutando in formulis [1], [4], [5], quantitates, quae referuntur ad \(B\), in eas, quae referuntur ad \(C\), nanciscemur formulas pro \(r'r',r'\cos\varphi',r'\sin\varphi'\). Tunc evolutio expressionis \(rr+r'r'-(q-q')^2-2r\cos\varphi\mathord{.}r'\cos\varphi'-2r\sin\varphi\mathord{.}r'\sin\varphi'\), quae fit \(=bb+cc-aa-2bc\cos A=2bc(\cos A^*-\cos A)\), combinata cum evolutione
